\documentclass[11pt]{article}
\usepackage[margin=0.85in]{geometry}
\usepackage{booktabs,longtable,array}
\usepackage{xcolor}
\usepackage{hyperref}
\usepackage{pdfpages}
\hypersetup{colorlinks=true,linkcolor=blue,urlcolor=blue}
\emergencystretch=2em

\title{ELMFIRE Verification Summary Report}
\author{ELMFIRE Verification and Validation Suite}
\date{\today}

\begin{document}
\maketitle

\begin{abstract}
This document aggregates the standalone reports for every runnable verification
case currently discovered under \texttt{cases/Verification}.  The opening
summary records environment and execution configuration and reports each case's
scientific decision from its case-local \texttt{outputs/metrics.json}; command
completion and PDF availability are not substituted for verification evidence.
\end{abstract}

\tableofcontents
\clearpage

% Generated by tools/generate_summary_reports.py; do not edit by hand.
\clearpage
\section{Verification summary}
\label{sec:verification-summary}
\subsection{Environment captured for this aggregate}
This environment was captured when the aggregate report was generated. It describes the report-build process and must not be interpreted as the execution environment of every case unless the cases were run and aggregated in the same job. ELMFIRE is not executed to obtain this information.

\small
\begin{longtable}{@{}>{\raggedright\arraybackslash}p{0.27\textwidth}>{\raggedright\arraybackslash}p{0.68\textwidth}@{}}
\toprule
Configuration & Recorded value \\
\midrule
\endfirsthead
\toprule
Configuration & Recorded value \\
\midrule
\endhead
\bottomrule
\endlastfoot
Captured at (UTC) & 2026-09-03T18:18:33+00:00 \\
Repository revision & 3304df6c77bb \\
Repository state & dirty \\
Host platform & macOS-15.7.7-arm64-arm-64bit \\
Python runtime & CPython 3.12.9 (/\allowbreak{}opt/\allowbreak{}homebrew/\allowbreak{}anaconda3/\allowbreak{}bin/\allowbreak{}python3) \\
Conda environment & not active \\
ELMFIRE command & elmfire \\
ELMFIRE resolved path & not found on PATH \\
ELMFIRE version declaration & not declared \\
ELMFIRE executable SHA-256 & not available \\
MPI launcher & not available \\
Slurm allocation & not a Slurm build job \\
Loaded modules & not recorded \\
Python packages & NumPy 1.26.4; \allowbreak{}Matplotlib 3.9.2; \allowbreak{}PyYAML 6.0.1; \allowbreak{}Rasterio 1.5.0; \allowbreak{}GDAL not installed; \allowbreak{}Pandas 2.2.2; \allowbreak{}GeoPandas not installed; \allowbreak{}Shapely not installed \\
\end{longtable}
\normalsize

\subsection{Case-declared execution configuration}
These values come only from each case's \texttt{case.yaml}. They summarize the declared executable command, namelist, MPI allocation, and namelist contract. A missing declaration is reported explicitly rather than inferred.

\scriptsize
\begin{longtable}{@{}>{\raggedright\arraybackslash}p{0.17\textwidth}>{\raggedright\arraybackslash}p{0.21\textwidth}>{\raggedright\arraybackslash}p{0.23\textwidth}>{\raggedright\arraybackslash}p{0.12\textwidth}>{\raggedright\arraybackslash}p{0.12\textwidth}@{}}
\toprule
Case & ELMFIRE command & Input namelist & MPI ranks & Schema \\
\midrule
\endfirsthead
\toprule
Case & ELMFIRE command & Input namelist & MPI ranks & Schema \\
\midrule
\endhead
\midrule
\multicolumn{5}{r}{Continued on next page} \\
\endfoot
\bottomrule
\endlastfoot
\hyperref[case-detail:case01-bet]{CASE01\_BET} & \$\{ELMFIRE\_BIN:-elmfire\} & elmfire.data.in & not declared & 1 \\
\hyperref[case-detail:case02-fgr]{CASE02\_FGR} & \$\{ELMFIRE\_BIN:-elmfire\} & elmfire.data.in & not declared & 1 \\
\hyperref[case-detail:case03-etc]{CASE03\_ETC} & \$\{ELMFIRE\_BIN:-elmfire\} & elmfire.data.in & not declared & 1 \\
\hyperref[case-detail:case04-sti]{CASE04\_STI} & \$\{ELMFIRE\_BIN:-elmfire\} & elmfire.data.in & not declared & 1 \\
\hyperref[case-detail:case05-tdw]{CASE05\_TDW} & \$\{ELMFIRE\_BIN:-elmfire\} & elmfire.data.in & not declared & 1 \\
\hyperref[case-detail:case06-f2d]{CASE06\_F2D} & \$\{ELMFIRE\_BIN:-elmfire\} & elmfire.data.in & not declared & 1 \\
\hyperref[case-detail:case07-ipc]{CASE07\_IPC} & \$\{ELMFIRE\_BIN:-elmfire\} & elmfire.data.in & not declared & 1 \\
\hyperref[case-detail:case08-fbc]{CASE08\_FBC} & \$\{ELMFIRE\_BIN:-elmfire\} & elmfire.data.in & not declared & 1 \\
\hyperref[case-detail:case09-mmd]{CASE09\_MMD} & \$\{ELMFIRE\_BIN:-elmfire\} & elmfire.data.in & not declared & 1 \\
\hyperref[case-detail:case10-wsr]{CASE10\_WSR} & \$\{ELMFIRE\_BIN:-elmfire\} & elmfire.data.in & not declared & 1 \\
\hyperref[case-detail:case11-wsc]{CASE11\_WSC} & \$\{ELMFIRE\_BIN:-elmfire\} & elmfire.data.in & not declared & 1 \\
\hyperref[case-detail:case12-trs]{CASE12\_TRS} & \$\{ELMFIRE\_BIN:-elmfire\} & elmfire.data.in & not declared & 1 \\
\hyperref[case-detail:case13-prm]{CASE13\_PRM} & not applicable & not applicable & not declared & 1 \\
\hyperref[case-detail:case14-srs]{CASE14\_SRS} & \$\{ELMFIRE\_BIN:-elmfire\} & elmfire.data.in & not declared & 1 \\
\hyperref[case-detail:case15-cro]{CASE15\_CRO} & \$\{ELMFIRE\_BIN:-elmfire\} & data/\allowbreak{}inputs/\allowbreak{}elmfire.data & not declared & 1 \\
\hyperref[case-detail:case16-nsp]{CASE16\_NSP} & \$\{ELMFIRE\_BIN:-elmfire\} & data/\allowbreak{}inputs/\allowbreak{}elmfire.data & not declared & 1 \\
\hyperref[case-detail:case17-scv]{CASE17\_SCV} & \$\{ELMFIRE\_BIN:-elmfire\} & data/\allowbreak{}128/\allowbreak{}inputs/\allowbreak{}elmfire.data & not declared & 1 \\
\hyperref[case-detail:case18-tcv]{CASE18\_TCV} & \$\{ELMFIRE\_BIN:-elmfire\} & data/\allowbreak{}inputs/\allowbreak{}elmfire\_180.data & not declared & 1 \\
\hyperref[case-detail:case19-wth]{CASE19\_WTH} & \$\{ELMFIRE\_BIN:-elmfire\} & elmfire.data.in & not declared & 1 \\
\hyperref[case-detail:case20-pig]{CASE20\_PIG} & \$\{ELMFIRE\_BIN:-elmfire\} & elmfire.data.in & not declared & 1 \\
\hyperref[case-detail:case21-wde]{CASE21\_WDE} & \$\{ELMFIRE\_BIN:-elmfire\} & elmfire.data.in & not declared & 1 \\
\hyperref[case-detail:case22-vsl]{CASE22\_VSL} & \$\{ELMFIRE\_BIN:-elmfire\} & elmfire.data.in & not declared & 1 \\
\hyperref[case-detail:case23-fqd]{CASE23\_FQD} & \$\{ELMFIRE\_BIN:-elmfire\} & elmfire.data.in & not declared & 1 \\
\hyperref[case-detail:case24-cxi]{CASE24\_CXI} & \$\{ELMFIRE\_BIN:-elmfire\} & elmfire.data.in & not declared & 1 \\
\hyperref[case-detail:case25-mqd]{CASE25\_MQD} & \$\{ELMFIRE\_BIN:-elmfire\} & elmfire.data.in & not declared & 1 \\
\hyperref[case-detail:case26-cnf]{CASE26\_CNF} & \$\{ELMFIRE\_BIN:-elmfire\} & variants/\allowbreak{}1mph/\allowbreak{}elmfire.data.in & not declared & 1 \\
\hyperref[case-detail:case27-fbt]{CASE27\_FBT} & \$\{ELMFIRE\_BIN:-elmfire\} & elmfire.data.in & not declared & 1 \\
\hyperref[case-detail:case28-ova]{CASE28\_OVA} & \$\{ELMFIRE\_BIN:-elmfire\} & elmfire.data.in & not declared & 1 \\
\hyperref[case-detail:case29-sup]{CASE29\_SUP} & \$\{ELMFIRE\_BIN:-elmfire\} & elmfire.data.in & not declared & 1 \\
\hyperref[case-detail:case30-mit]{CASE30\_MIT} & \$\{ELMFIRE\_BIN:-elmfire\} & elmfire.data.in & not declared & 1 \\
\hyperref[case-detail:case31-pft]{CASE31\_PFT} & \$\{ELMFIRE\_BIN:-elmfire\} & elmfire.data.in & not declared & 1 \\
\hyperref[case-detail:case32-rcv]{CASE32\_RCV} & \$\{ELMFIRE\_BIN:-elmfire\} & elmfire.data.in & not declared & 1 \\
\hyperref[case-detail:case33-sdo]{CASE33\_SDO} & \$\{ELMFIRE\_BIN:-elmfire\} & elmfire.data.in & not declared & 1 \\
\hyperref[case-detail:case34-fms]{CASE34\_FMS} & \$\{ELMFIRE\_BIN:-elmfire\} & elmfire.data.in & not declared & 1 \\
\hyperref[case-detail:case35-wss]{CASE35\_WSS} & \$\{ELMFIRE\_BIN:-elmfire\} & elmfire.data.in & not declared & 1 \\
\hyperref[case-detail:case36-sls]{CASE36\_SLS} & \$\{ELMFIRE\_BIN:-elmfire\} & elmfire.data.in & not declared & 1 \\
\hyperref[case-detail:case37-dms]{CASE37\_DMS} & \$\{ELMFIRE\_BIN:-elmfire\} & elmfire.data.in & not declared & 1 \\
\hyperref[case-detail:case38-lms]{CASE38\_LMS} & \$\{ELMFIRE\_BIN:-elmfire\} & elmfire.data.in & not declared & 1 \\
\hyperref[case-detail:case39-dhc]{CASE39\_DHC} & \$\{ELMFIRE\_BIN:-elmfire\} & elmfire.data.in & not declared & 1 \\
\hyperref[case-detail:case40-wsd]{CASE40\_WSD} & \$\{ELMFIRE\_BIN:-elmfire\} & elmfire.data.in & not declared & 1 \\
\hyperref[case-detail:case41-cfp]{CASE41\_CFP} & \$\{ELMFIRE\_BIN:-elmfire\} & elmfire.data.in & not declared & 1 \\
\hyperref[case-detail:case42-waf]{CASE42\_WAF} & \$\{ELMFIRE\_BIN:-elmfire\} & elmfire.data.in & not declared & 1 \\
\end{longtable}
\normalsize

\subsection{Scientific decision summary}
This table is generated from each case's case-local \texttt{outputs/metrics.json}. A successful run script or report build is not treated as scientific evidence of PASS. Missing, incomplete, blocked, or non-decisional results are reported as NOT EVALUABLE.

\textbf{Total cases:} 42\quad
\textbf{PASS:} 30\quad
\textbf{FAIL:} 8\quad
\textbf{CHARACTERIZED:} 0\quad
\textbf{NOT EVALUABLE:} 4

\small
\setlength{\LTpre}{0.8em}
\setlength{\LTpost}{0pt}
\begin{longtable}{@{}>{\raggedright\arraybackslash}p{0.15\textwidth}>{\raggedright\arraybackslash}p{0.16\textwidth}>{\raggedright\arraybackslash}p{0.38\textwidth}>{\raggedright\arraybackslash}p{0.22\textwidth}@{}}
\toprule
Case & Category & Title & Decision \\
\midrule
\endfirsthead
\toprule
Case & Category & Title & Decision \\
\midrule
\endhead
\midrule
\multicolumn{4}{r}{Continued on next page} \\
\endfoot
\bottomrule
\endlastfoot
\hyperref[case-detail:case01-bet]{CASE01\_BET} & coupling tests & \hyperref[case-detail:case01-bet]{Biomass Emission-Time Resolution Independence} & \mbox{\textcolor{green!45!black}{\textbf{PASS}}} \\
\hyperref[case-detail:case02-fgr]{CASE02\_FGR} & coupling tests & \hyperref[case-detail:case02-fgr]{Finite-Duration Firebrand Generation} & \mbox{\textcolor{green!45!black}{\textbf{PASS}}} \\
\hyperref[case-detail:case03-etc]{CASE03\_ETC} & coupling tests & \hyperref[case-detail:case03-etc]{Eulerian Firebrand Transport Convergence} & \mbox{\textcolor{red!75!black}{\textbf{FAIL}}} \\
\hyperref[case-detail:case04-sti]{CASE04\_STI} & coupling tests & \hyperref[case-detail:case04-sti]{Steady Transport and Immediate Ignition} & \mbox{\textcolor{green!45!black}{\textbf{PASS}}} \\
\hyperref[case-detail:case05-tdw]{CASE05\_TDW} & coupling tests & \hyperref[case-detail:case05-tdw]{Firebrand Transport in Time-Dependent Wind} & \mbox{\textcolor{green!45!black}{\textbf{PASS}}} \\
\hyperref[case-detail:case06-f2d]{CASE06\_F2D} & coupling tests & \hyperref[case-detail:case06-f2d]{Two-Dimensional Firebrand-Driven Spread} & \mbox{\textcolor{green!45!black}{\textbf{PASS}}} \\
\hyperref[case-detail:case07-ipc]{CASE07\_IPC} & coupling tests & \hyperref[case-detail:case07-ipc]{Ignition Probability and SFT Convergence} & \mbox{\textcolor{red!75!black}{\textbf{FAIL}}} \\
\hyperref[case-detail:case08-fbc]{CASE08\_FBC} & coupling tests & \hyperref[case-detail:case08-fbc]{Deposited Firebrand Consumption} & \mbox{\textcolor{red!75!black}{\textbf{FAIL}}} \\
\hyperref[case-detail:case09-mmd]{CASE09\_MMD} & coupling tests & \hyperref[case-detail:case09-mmd]{Mixed-Mode Ignition-Delay Effects} & \mbox{\textcolor{green!45!black}{\textbf{PASS}}} \\
\hyperref[case-detail:case10-wsr]{CASE10\_WSR} & coupling tests & \hyperref[case-detail:case10-wsr]{WUI Structural Resolution Sensitivity} & \mbox{\textcolor{orange!75!black}{\textbf{NOT EVALUABLE}}} \\
\hyperref[case-detail:case11-wsc]{CASE11\_WSC} & coupling tests & \hyperref[case-detail:case11-wsc]{WUI Surface-Coupling Synchronization} & \mbox{\textcolor{orange!75!black}{\textbf{NOT EVALUABLE}}} \\
\hyperref[case-detail:case12-trs]{CASE12\_TRS} & coupling tests & \hyperref[case-detail:case12-trs]{WUI Temporal-Resolution Sensitivity} & \mbox{\textcolor{red!75!black}{\textbf{FAIL}}} \\
\hyperref[case-detail:case13-prm]{CASE13\_PRM} & coupling tests & \hyperref[case-detail:case13-prm]{WUI Parametric Response Maps} & \mbox{\textcolor{orange!75!black}{\textbf{NOT EVALUABLE}}} \\
\hyperref[case-detail:case14-srs]{CASE14\_SRS} & coupling tests & \hyperref[case-detail:case14-srs]{Two-Dimensional Spatial-Resolution Sensitivity} & \mbox{\textcolor{orange!75!black}{\textbf{NOT EVALUABLE}}} \\
\hyperref[case-detail:case15-cro]{CASE15\_CRO} & unit tests & \hyperref[case-detail:case15-cro]{Case 15: Constant rate of spread} & \mbox{\textcolor{green!45!black}{\textbf{PASS}}} \\
\hyperref[case-detail:case16-nsp]{CASE16\_NSP} & unit tests & \hyperref[case-detail:case16-nsp]{Case 16: No-spread limiting behavior} & \mbox{\textcolor{green!45!black}{\textbf{PASS}}} \\
\hyperref[case-detail:case17-scv]{CASE17\_SCV} & unit tests & \hyperref[case-detail:case17-scv]{Case 17: Spatial convergence} & \mbox{\textcolor{red!75!black}{\textbf{FAIL}}} \\
\hyperref[case-detail:case18-tcv]{CASE18\_TCV} & unit tests & \hyperref[case-detail:case18-tcv]{Case 18: Temporal convergence} & \mbox{\textcolor{red!75!black}{\textbf{FAIL}}} \\
\hyperref[case-detail:case19-wth]{CASE19\_WTH} & coupling tests & \hyperref[case-detail:case19-wth]{Case 19: Transient WUI heat flux} & \mbox{\textcolor{green!45!black}{\textbf{PASS}}} \\
\hyperref[case-detail:case20-pig]{CASE20\_PIG} & coupling tests & \hyperref[case-detail:case20-pig]{Case 20: Point ignition and isotropic spread} & \mbox{\textcolor{green!45!black}{\textbf{PASS}}} \\
\hyperref[case-detail:case21-wde]{CASE21\_WDE} & coupling tests & \hyperref[case-detail:case21-wde]{Case 21: Wind-driven elliptical spread} & \mbox{\textcolor{green!45!black}{\textbf{PASS}}} \\
\hyperref[case-detail:case22-vsl]{CASE22\_VSL} & coupling tests & \hyperref[case-detail:case22-vsl]{Case 22: Spread across opposing valley slopes} & \mbox{\textcolor{green!45!black}{\textbf{PASS}}} \\
\hyperref[case-detail:case23-fqd]{CASE23\_FQD} & coupling tests & \hyperref[case-detail:case23-fqd]{Case 23: Spatially varying fuel quadrants} & \mbox{\textcolor{green!45!black}{\textbf{PASS}}} \\
\hyperref[case-detail:case24-cxi]{CASE24\_CXI} & coupling tests & \hyperref[case-detail:case24-cxi]{Case 24: Combined fuel, slope, and wind interactions} & \mbox{\textcolor{green!45!black}{\textbf{PASS}}} \\
\hyperref[case-detail:case25-mqd]{CASE25\_MQD} & coupling tests & \hyperref[case-detail:case25-mqd]{Case 25: Spatially varying moisture quadrants} & \mbox{\textcolor{green!45!black}{\textbf{PASS}}} \\
\hyperref[case-detail:case26-cnf]{CASE26\_CNF} & coupling tests & \hyperref[case-detail:case26-cnf]{Case 26: Surface, passive-crown, and active-crown fire} & \mbox{\textcolor{green!45!black}{\textbf{PASS}}} \\
\hyperref[case-detail:case27-fbt]{CASE27\_FBT} & coupling tests & \hyperref[case-detail:case27-fbt]{Case 27: Lagrangian firebrand generation and transport} & \mbox{\textcolor{red!75!black}{\textbf{FAIL}}} \\
\hyperref[case-detail:case28-ova]{CASE28\_OVA} & coupling tests & \hyperref[case-detail:case28-ova]{Case 28: Overnight spread-rate adjustment} & \mbox{\textcolor{green!45!black}{\textbf{PASS}}} \\
\hyperref[case-detail:case29-sup]{CASE29\_SUP} & coupling tests & \hyperref[case-detail:case29-sup]{Case 29: Initial- and extended-attack suppression} & \mbox{\textcolor{red!75!black}{\textbf{FAIL}}} \\
\hyperref[case-detail:case30-mit]{CASE30\_MIT} & coupling tests & \hyperref[case-detail:case30-mit]{Case 30: Multiple-ignition topology} & \mbox{\textcolor{green!45!black}{\textbf{PASS}}} \\
\hyperref[case-detail:case31-pft]{CASE31\_PFT} & coupling tests & \hyperref[case-detail:case31-pft]{Case 31: Planar-front transport} & \mbox{\textcolor{green!45!black}{\textbf{PASS}}} \\
\hyperref[case-detail:case32-rcv]{CASE32\_RCV} & coupling tests & \hyperref[case-detail:case32-rcv]{Case 32: Raster-orientation covariance} & \mbox{\textcolor{green!45!black}{\textbf{PASS}}} \\
\hyperref[case-detail:case33-sdo]{CASE33\_SDO} & coupling tests & \hyperref[case-detail:case33-sdo]{Case 33: Slope-driven directional spread} & \mbox{\textcolor{green!45!black}{\textbf{PASS}}} \\
\hyperref[case-detail:case34-fms]{CASE34\_FMS} & unit tests & \hyperref[case-detail:case34-fms]{Case 34: Standard fuel-model Rothermel sweep} & \mbox{\textcolor{green!45!black}{\textbf{PASS}}} \\
\hyperref[case-detail:case35-wss]{CASE35\_WSS} & unit tests & \hyperref[case-detail:case35-wss]{Case 35: Midflame wind-speed Rothermel sweep} & \mbox{\textcolor{green!45!black}{\textbf{PASS}}} \\
\hyperref[case-detail:case36-sls]{CASE36\_SLS} & unit tests & \hyperref[case-detail:case36-sls]{Case 36: Slope-factor Rothermel sweep} & \mbox{\textcolor{green!45!black}{\textbf{PASS}}} \\
\hyperref[case-detail:case37-dms]{CASE37\_DMS} & unit tests & \hyperref[case-detail:case37-dms]{Case 37: Dead-fuel moisture and extinction sweep} & \mbox{\textcolor{green!45!black}{\textbf{PASS}}} \\
\hyperref[case-detail:case38-lms]{CASE38\_LMS} & unit tests & \hyperref[case-detail:case38-lms]{Case 38: Static live-fuel moisture sweep} & \mbox{\textcolor{green!45!black}{\textbf{PASS}}} \\
\hyperref[case-detail:case39-dhc]{CASE39\_DHC} & unit tests & \hyperref[case-detail:case39-dhc]{Case 39: Dynamic herbaceous curing sweep} & \mbox{\textcolor{green!45!black}{\textbf{PASS}}} \\
\hyperref[case-detail:case40-wsd]{CASE40\_WSD} & coupling tests & \hyperref[case-detail:case40-wsd]{Case 40: Wind-slope vector and direction sweep} & \mbox{\textcolor{green!45!black}{\textbf{PASS}}} \\
\hyperref[case-detail:case41-cfp]{CASE41\_CFP} & unit tests & \hyperref[case-detail:case41-cfp]{Case 41: Custom fuel-parameter Rothermel sweep} & \mbox{\textcolor{green!45!black}{\textbf{PASS}}} \\
\hyperref[case-detail:case42-waf]{CASE42\_WAF} & coupling tests & \hyperref[case-detail:case42-waf]{Case 42: Wind-adjustment-factor coupling sweep} & \mbox{\textcolor{green!45!black}{\textbf{PASS}}} \\
\end{longtable}
\normalsize

\begin{samepage}
\paragraph{Decision vocabulary.}
PASS and FAIL are case-defined scientific decisions. CHARACTERIZED denotes descriptive validation metrics without a justified acceptance threshold. NOT EVALUABLE denotes missing, unrun, incomplete, unsupported, or otherwise non-decisional evidence.
\end{samepage}


\section{Scope and decision policy}

The report order follows the stable numerical CASE identifier.  Each standalone
case PDF is included intact to preserve its equations, figures, acceptance
criteria, and limitations.  PASS and FAIL retain the decision made by the
case-local postprocessor.  A missing result, an unrun or incomplete workflow,
or an unsupported capability is shown as NOT EVALUABLE.  The linked summary at
the beginning records the aggregate-build environment, case-declared execution
configuration, and scientific decisions before the complete case reports.

\section{How to read this report}

Each entry is the complete standalone PDF for one verification case.  Cases
follow the stable order in
\path{cases/Verification/CASE_REGISTRY.md}, and each report retains its own
scientific context, equations, metrics, citations, and decision.  Aggregating
the compiled PDFs keeps case-local macros, labels, bibliographies, and packages
isolated from one another.  After standalone reports change, rebuild this guide
with \texttt{make verification-report} or rebuild both suite reports with
\texttt{make reports}.

\input{generated/verification_cases.tex}

\end{document}
